% Part: lambda-calculus
% Chapter: introduction
% Section: beta

\documentclass[../../../include/open-logic-section]{subfiles}

\begin{document}

\olfileid{lam}{int}{bet}
\olsection{$\beta$-راکمول}

کله چې $(\lambd[m][(\lambd[y][y]) m])$ ووينو، طبيعي ده چې
ګومان وکړو له $\lambd[m][m]$ سره څه اړيکه لري؛ يعنې دويم ترم
بايد د لومړي د «اسانولو» پايله وي۔ د \emph{$\beta$-راکمول}
مفهوم دا انګېرنه په رسمي ډول رانغاړي۔

\begin{defn}[$\beta$-انقباض، $\bredone$] \ollabel{defn:betacontr}
  \emph{$\beta$-انقباض} ($\bredone$) پر ترمونو تر ټولو کوچنۍ
  سازګاره اړيکه ده چې لاندې شرط پوره کوي:
  \[
    (\lambd[x][N])Q \bredone \Subst{N}{Q}{x} 
  \]
  که $P \bredone Q$ وي، وايو چې $P$، $Q$ ته
  \emph{$\beta$-منقبض} شوے دے۔ د $(\lambd[x][N])Q$
  په بڼه ترم ته \emph{رېډکس} وايي۔
\end{defn}

\begin{prob} \ollabel{prob:def}
  د $\beta$-انقباض معادل استقرايي تعريفونه په څرګنده وليکئ،
  لکه څنګه چې مو د تړلي متغير د نوم بدلون دپاره په
  \olref[lam][syn][alp]{defn:aconvone} کښې وکړل۔
\end{prob}
  
\begin{defn}[$\beta$-راکمول، $\bred$] \ollabel{defn:betared}
  \emph{$\beta$-راکمول} ($\bred$) پر ترمونو تر ټولو کوچنۍ
  انعکاسي او انتقالي اړيکه ده چې $\bredone$ پکې شامله وي۔
  که $P \bred Q$ وي، وايو چې $P$، $Q$ ته
  $\beta$-راکم شوے دے۔
\end{defn}

کله چې مطلب له متنه څرګند وي، د $\bredone$ پر ځاے
$\redone$ او د $\bred$ پر ځاے $\red$ ليکو۔

په غيررسمي وينا، $M \bred N$ هله او يوازې هله وي چې
$M$ د $\beta$-انقباض په صفر يا څو ګامونو $N$ ته
بدلېدے شي۔

\begin{defn}[$\beta$-عادي]
هغه ترم چې نور پرې $\beta$-انقباض نۀ شي کېدے،
\emph{$\beta$-عادي} بلل کېږي۔
\end{defn}

که $M \bred N$ او $N$، $\beta$-عادي وي، نو $N$
د~$M$ \emph{عادي بڼه} بولو۔ دا پوښتنه پيدا کېږي چې
د يوه ترم عادي بڼه يکتا ده که نه؛ ځواب هو دے، لکه
وروسته به ئې ووينو۔
\textbf{د ژباړې د نسخې سپيناوی (OLLAM-032):} دا د
يکتا والي خبره د عادي بڼې د موجوديت په شرط ده؛
هر ترم عادي بڼې ته نۀ رسېږي۔ د سرچينې يوه
وروستۍ برخه همدغه شرط په څرګنده وايي۔

راځئ څو بېلګې وګورو۔
\begin{enumerate}
\item لرو:
\begin{align*}
(\lambd[x][xxy]) \lambd[z][z] & \redone (\lambd[z][z])(\lambd[z][z]) y \\
& \redone (\lambd[z][z]) y \\
& \redone y
\end{align*}
\item د يوه ترم «اسانول» په حقيقت کښې هغه لا پېچلے کولے شي:
\begin{align*}
(\lambd[x][xxy])(\lambd[x][xxy]) & \redone (\lambd[x][xxy])(\lambd[x][xxy])y \\
& \redone (\lambd[x][xxy])(\lambd[x][xxy])yy \\
& \redone \dots
\end{align*}
\item دا کار ترم هماغسې بې بدله هم پرېښودے شي:
\[
(\lambd[x][xx])(\lambd[x][xx]) \redone (\lambd[x][xx])(\lambd[x][xx])
\]
\item ځينې ترمونه په يوه نه، بلکې په څو لارو هم راکمېږي؛
  د بېلګې په توګه،
\[
(\lambd[x][(\lambd[y][yx]) z]) v \redone (\lambd[y][yv]) z
\]
  د تر ټولو دباندني تطبيق د انقباض له لارې؛ او
\[
(\lambd[x][(\lambd[y][yx]) z]) v \redone (\lambd[x][zx]) v
\]
  د تر ټولو دننني تطبيق د انقباض له لارې۔ خو پام وکړئ
  چې په دې حالت کښې دواړه پايلې بيا هماغه يوه ترم،
  $zv$، ته راکمېږي۔
\end{enumerate}

د وروستۍ بېلګې ګډه پايله تصادف نۀ دے؛ دا د لامبډا
حساب يو ژور او مهم خاصيت ښيي چې د چرچ--روسر خاصيت
بلل کېږي۔

\begin{digress}
  په عمومي ډول يو ترم د $\beta$-راکمول څو لارې لرلے
  شي؛ له همدې امله د راکمول بېلابېلې تګلارې جوړې
  شوې دي۔ په هغو کښې تر ټولو عامه \emph{طبيعي تګلاره}
  ده۔ طبيعي تګلاره تل هغه \emph{تر ټولو کيڼ} رېډکس
  منقبضوي چې ځای ئې د ترم په ليکلي بڼه کښې د هغه
  د پيل له مخې ټاکل کېږي۔ د دې تګلارې ګټور خاصيت
  دا دے: که يو ترم په کومې تګلارې عادي بڼې ته
  راکمېدے شي، نو په طبيعي تګلارې هم هغې ته
  راکمېدے شي۔ له دې وروسته به، تر څو بل ډول
  څرګند نۀ شي، همدا طبيعي تګلاره کاروو۔
\end{digress}

\begin{defn}[$\beta$-معادلتوب، $\equal$]
  \emph{$\beta$-معادلتوب} ($\equal$) هغه اړيکه ده چې
  په استقرايي ډول داسې تعريفېږي:
  \begin{enumerate}
  \item $M \equal M$.
  \item که $M \equal N$ وي، نو $N \equal M$.
  \item که $M \equal N$ او $N \equal O$ وي، نو $M \equal O$.
  \item که $M \equal N$ وي، نو $PM \equal PN$.
  \item که $M \equal N$ وي، نو $MQ \equal NQ$.
  \item که $M \equal N$ وي، نو $\lambd[x][M] \equal \lambd[x][N]$.
  \item $(\lambd[x][N])Q \equal \Subst{N}{Q}{x}$.
  \end{enumerate}
\end{defn}

لومړۍ درې قاعدې دا اړيکه د معادلتوب اړيکه ګرځوي؛
راتلونکې درې ئې سازګاره کوي؛ او وروستۍ قاعده
$\beta$-انقباض پکې شاملوي۔

په غيررسمي وينا، دوه ترمونه هله او يوازې هله
$\beta$-معادل دي چې يو ئې د $\beta$-انقباض يا د
معکوس $\beta$-انقباض په صفر يا څو ګامونو بل ته بدل شي۔
د $\beta$-انقباض معکوس داسې تعريفېږي چې $M$،
$N$ ته په معکوس $\beta$-انقباض هله او يوازې هله
اوړي چې $N$، $M$ ته $\beta$-منقبضېږي۔

له پورته قواعدو سربېره به دې اړيکې ته نور قواعد
هم ورزياتوو او غځېدلې معادلتوبي اړيکه به
$\equal[X]$ ليکو؛ دلته $X$ هغه ورزياته شوې قاعده ده۔

\end{document}
